\chapter{Discourse, Register, and Spoken Iridian}
\label{ch:discourse}\label{ch:spoken}\label{ch:semantics}

% ============================================================
% This chapter brings together three domains that belong to a single
% descriptive continuum: how Iridian organises information across clauses
% and texts; how register and social context shape grammatical choices;
% and how the spoken and colloquial variety departs from the written
% standard. The three domains interact: aspect choice as a narrative
% device is also a discourse-register choice; many spoken morphemes are
% formally integrated into the same inflectional machinery that the rest
% of the grammar describes; and politeness strategies depend on the same
% anchoring and topic mechanisms analysed in earlier chapters.
% ============================================================

\section{Introduction}\label{sec:discourse-intro}

Iridian maintains reference across clauses through overt nominals,
 demonstratives, reflexives, and recoverable null arguments. Aspect and voice
 determine the grammatical pivot of each predicate, while discourse determines
 which participant remains the current subject matter. The two often
 coincide, although linked agent topics and external resumptive topics permit
 their separation (Chapter~\ref{ch:infstr}). Register and speaker stance
 further condition the choice of particles and the use of colloquial forms.

%
% ================================================================
%  PART I: DISCOURSE STRUCTURE
% ================================================================
%

\section{Sentence chains as discourse units}\label{sec:discourse-chains}

% ============================================================
% Iridian discourse is organised in "sentence chains": sequences of
% converb clauses capped by a finite verb. Each sentence chain functions
% as a single discourse unit. The chain has:
%   - Converb clauses = background, supporting, or subordinate events
%   - Finite verb = the foregrounded, topically relevant event
% This parallels the foreground/background distinction found in many
% verb-final or SOV languages with converb systems.
% ============================================================

\section{Topic continuity and topic shift}\label{sec:topic-continuity}


\subsection{Reference management without morphological cross-reference}
\label{sec:ref-management}

Iridian manages reference tracking without any pronominal clitic series on the
verb or noun. The devices available are: (a)~overt demonstratives, which encode
the proximate/obviative distinction and serve as the primary third-person
referential series; (b)~overt first- and second-person pronouns, employed when
required by contrast, emphasis, disambiguation, or low discourse recoverability;
(c)~the reflexive anaphor, which forces local coreference within the clausal
domain; and (d)~zero anaphora, the unmarked strategy in contexts of stable topic
continuity, where the intended argument is recoverable
without any overt expression.

Pro-drop in Iridian is accordingly a purely discourse-driven null-argument
system. The verb does not morphologically cross-reference its arguments through
a clitic or agreement layer: argument drop is licensed entirely by discourse
recoverability and topic continuity, and the burden of disambiguation falls on
syntax and pragmatics rather than on right-edge person morphology.

\section{Anchoring in discourse context}\label{sec:discourse-definiteness}

% The anchoring suffix -to on a noun signals that the referent is
% identifiable to the hearer. In discourse, first mention is typically
% bare (no anchoring marker); subsequent mention may use the marker,
% a demonstrative, or drop the noun entirely (zero anaphora).
% The interaction of anchoring, topic, and absolutive alignment
% is the key thread here.

\subsection{The neutral discourse centre in direct-case position}
\label{sec:disc-def}

The suffixal anchor \ird{-to}/\ird{-te} maintains an identifiable nominal
 referent, whether that referent is a topic, a backgrounded participant, or
 part of the asserted information. Direct-case patients and transitive-case
 agents have the same shared-accessibility contrast; their case does not add
 contrastive or resumptive meaning to the anchor
 (§\,\ref{sec:nm-def-nondirect}). A linked topic can independently promote
 an anchored agent to the subject matter of a perfective clause, while the
 patient retains direct case (§\,\ref{sec:tc-marked-topics}).

A continuing aboutness topic supplies a strong contextual antecedent for a
 recoverable null argument. Where no such topic has been independently
 established, the preceding direct-case pivot provides the neutral fallback.
 These are defeasible discourse preferences: the predicate's interpretation,
 competing referents, and overt reintroduction constrain reference. They do
 not replace the same-pivot relation in converb chains
 (§\,\ref{sec:switch-reference}).

\section{Aspect choice as narrative strategy}\label{sec:aspect-narrative}

Perfective and retrospective clauses commonly advance a sequence of bounded
 events, whereas imperfective clauses can describe ongoing activities or
 background states. Their different case frames also change the default
 pivot candidate for discourse continuation. That association does not force
 a change of topic: a perfective agent can remain the subject matter of an
 independently topicalised predication, and an imperfective description can
 introduce a new referent (§\,\ref{sec:tc-marked-topics}).

The antipassive supports agent continuity by retaining the agent as pivot
while removing the patient from core argument structure. It presents an
event chiefly as part of the agent's activity, allowing that agent to remain
topical without establishing or tracking a core patient. It is especially natural when
the patient is generic, incorporated into the event type, inferable from
context, or simply of no further narrative consequence; for the morphological
details see Chapter~\ref{ch:verbs}.

The causative introduces a causer and presents the event as one of
inducement, delegated action, permission, or coercion. In narrative, this
additional participant can foreground authority, distribute responsibility,
or distinguish the discourse topic from the participant who directly
performs the action.
In perfective contexts this often yields a strongly interventionist reading; in
retrospective contexts it more readily supports permissive, indirect, or
weak-control interpretations. The formal argument structure of the causative is
described in Chapter~\ref{ch:verbs}. Its additional argument supplies another
potential centre for discourse continuation, whose selection as topic remains
independent of the causative morphology.

\subsection{Mini-texts: continuity and re-centering}

The following constructed sequences illustrate topic continuity where the
 topic coincides with the pivot, and a subsequent shift supported by overt
 reintroduction.

\pex
\a
\begingl
\gla Marek travě piaštime.//
\glb Marek.\Dir{} bread.\Ind{} eat-\Prog{}//
\glft \trsl{Marek is eating bread.}//
\endgl
\a
\begingl
\gla Zmacime.//
\glb run-\Prog{}//
\glft \trsl{[He] is running.} \lingcontext{agent continuity in the nominative domain}//
\endgl
\xe

\noindent Here the null topic of the second clause is naturally interpreted as
Marek, who is both the current aboutness topic and the direct-case pivot.
The predicate \ird{zmac-} also favours an animate participant, so the
inanimate bread is not an equivalent antecedent competitor.

\pex
\a
\begingl
\gla Tak tóm Marku koupnik.//
\glb \Dem{}.\Prox{}.\Dir{} book.\Dir{} Marek.\Trs{} buy-\Trs{}.\Pf{}//
\glft \trsl{This book, Marek bought.}//
\endgl
\a
\begingl
\gla By koupnik?//
\glb \textsc{q} buy-\Trs{}.\Pf{}//
\glft \trsl{Did [he] buy [it]?}
 \lingcontext{confirmation question continuing the same buyer and book}//
\endgl
\xe

\noindent In an exchange about the book, the confirmation question continues
 the patient topic while leaving the buyer unexpressed as well. The same
 perfective case frame holds in both clauses, and the discourse recovers each
 participant in the role selected by the verb. Since both referents remain
 accessible, the sequence illustrates continuity without isolating a special
 antecedent preference for the pivot. Overt reintroduction can instead make
 an agent the subject matter of the next clause, as in the following sequence.

\pex
\a
\begingl
\gla Travato Marku piaštnik.//
\glb bread.\Dir{}.\Def{} Marek.\Trs{} eat-\textsc{trs}-\Pf{}//
\glft \trsl{Marek ate the bread.}//
\endgl
\a
\begingl
\gla Marek travě piaštime.//
\glb Marek.\Dir{} bread.\Ind{} eat-\Prog{}//
\glft \trsl{Marek is eating bread.} \lingcontext{re-centering from patient-topic narration to agent-topic description}//
\endgl
\xe

\noindent The first clause can describe what happened to the bread, while
 the second explicitly reintroduces Marek as the participant whose current
 activity is described. The aspectual alternation and overt nominal jointly
 support that shift; neither variable has been held constant.

The established agent-topic clause \ird{Markuto tak tóm koupnik}
 (\ref{ex:tc-marked-agent-top}) separates the relations more directly:
 Marek is the aboutness topic and the book the grammatical pivot. A later null
 agent can be recovered as Marek when the discourse continues his activities,
 although a null patient in a matching buying clause can instead recover the
 book. The intended referent depends on its role in that clause, rather than
 on a mechanical instruction to copy the last direct-case expression.

The relative preference for a non-pivot topic and a competing pivot remains
 unresolved when both referents are equally plausible antecedents. The
 preceding sequences differ in animacy, predicate choice, or overt expression;
 they establish recoverability and overt shift, while leaving that competition
 under fixed aspect, case, and anchoring undetermined. The distinction from converb continuity
 remains grammatical: its same-pivot dependency is defined over the two
 predicates' pivots rather than over whichever participant the discourse is
 about (§\,\ref{sec:switch-reference}).

\section{Tail-head linkage}\label{sec:tail-head}

% Tail-head linkage is a discourse-structuring device where the final
% clause of one sentence chain is repeated (in whole or in abbreviated
% form) at the start of the next chain. This is attested in the sample
% texts and is a common device in narrative Iridian.

\section{Discourse particles}\label{sec:discourse-particles}\label{sec:discourse-word-order}


\subsection{Clause-final discourse particles}\label{sec:disc-final-particles}

Iridian permits a small class of clause-final discourse particles, placed after
the final prosodic constituent of the clause, most typically after the finite
verbal complex. These particles do not belong to the ordinary verbal
inflectional template and do not affect argument structure, alignment, or
truth-conditional content. Their function is instead pragmatic: they mark the
speaker's stance toward the proposition and toward the addressee at the level of
the utterance as a whole.

These particles stand at the utterance edge, outside the verbal stress domain.
In weak uses they may lean on the final prosodic word of the preceding clause;
emphatic uses permit independent prominence. Neither attachment nor emphasis
relocates lexical stress within that host (§\,\ref{sec:stress}). Semantically, they express such
values as appeal to shared knowledge, assertive presentation, corrective
insistence, or emotive intensification.

The principal clause-final particles are the following.

\paragraph{\ird{len}} The particle \ird{len} is confirmatory and shared-ground
oriented. It presents the proposition as accessible, recognisable, or available
for the addressee's ratification, with values such as `you know', `as we both
know', or `is that not so'. It is the least forceful member of the set and is
well suited to mild prompting, reassurance, and alignment of speaker and
addressee.

\paragraph{\ird{prav}} The particle \ird{prav} is assertive and presentational.
It marks the utterance as one the speaker puts forward with particular
commitment, roughly `indeed', `truly', or `mark this'. As a discourse particle
rather than a lexical predicate, it strengthens the assertion without taking an
argument of its own. In
clause-final position it signals strengthened assertion without necessarily
implying contrast.

\paragraph{\ird{však}} The particle \ird{však} is corrective or contrastive. It
indicates that the utterance stands against an assumption, implication, or line
of expectation in the discourse, with values such as `after all', `but in fact',
`mind you', or `contrary to that'. It is especially natural in rebuttal,
clarification, argumentative discourse, and narrative correction.

\paragraph{\ird{blěc}} The particle \ird{blěc} is a marked colloquial and
low-register emotive particle. Unlike \ird{len}, \ird{prav}, and \ird{však}, it
does not belong to the neutral core discourse system, but adds rough speaker
attitude such as irritation, incredulity, impatient insistence, or coarse
emphasis. Its contribution is primarily expressive and register-sensitive. In
formal or careful style it is normally avoided; in colloquial speech it may
reinforce commands, corrections, or exasperated assertions.

Taken together, these particles form a small utterance-final system in which
\ird{len} marks common-ground appeal, \ird{prav} assertive presentation,
\ird{však} corrective contrast, and \ird{blěc} emotive or vulgar
intensification. Further examples and distributional restrictions are discussed
below.

%
% ================================================================
%  PART II: REGISTER AND PRAGMATICS
% ================================================================
%

\section{Register and formality}\label{sec:register}\label{sec:register-discourse}


\section{Politeness and the T-V distinction}\label{sec:politeness}


\section{Honorifics}\label{sec:honorifics}


\section{Phatic expressions}\label{sec:phatic}

% [Greetings, leave-takings, thanks, apologies — cross-reference with
%  §\,\ref{sec:fixed-converbs} in Chapter~\ref{ch:complex}]

\section{Idiomatic expressions}\label{sec:idioms}


%
% ================================================================
%  PART III: SPOKEN AND COLLOQUIAL IRIDIAN
% ================================================================
%

\section{Spoken Iridian: an overview}\label{sec:spoken-overview}

This section sketches a small part of the spoken and urban colloquial register
of present-day Iridian. The written standard remains the reference variety of
this grammar, but it would be misleading to treat everyday spoken usage as
merely a noisier version of that standard. In ordinary conversation, speakers
make frequent use of reduced connectives, clause-peripheral hedges, semantically
bleached nouns, and other high-frequency colloquial forms whose contribution is
partly lexical and partly interactional.

Colloquial \ird{nu}, beside sequential \ird{nesu}, and the intentional use
of the purposive converb illustrate the interaction between lexical reduction
and discourse organisation (§§\,\ref{sec:seq-copula},
\ref{sec:purp-intentional}). The morphemes described below similarly combine
ordinary inflection with lexical or interactional meanings restricted to spoken
registers.

\section{Common colloquial morphemes}\label{sec:spoken-morphemes}

\subsection{The connective \ird{nu}}\label{sec:spoken-nu}

The colloquial connective \ird{nu} has already been introduced in
Section~\ref{sec:seq-copula} of Chapter~\ref{ch:complex}. Synchronically it is
best treated as the reduced spoken continuation of sequential \ird{nesu}, but
its discourse function extends beyond sequential clause-linking. In
rapid speech it often serves as a light connective between clauses, turns, or
discourse steps, with meanings ranging from loose \trsl{and then} to weaker
\trsl{so}, \trsl{well then}, or simple transition management.

Unlike the preverbal particles of Chapter~\ref{ch:function-words},
\ird{nu} occupies the edge of a clause or utterance, where it marks a
transition between discourse contributions.

\subsection{\ird{Žema} and reduced \ird{žma}}\label{sec:spoken-zema}

Among colloquial noun-like items, \ird{žema} is especially revealing. In
present-day spoken Iridian it is best analysed as a semantically bleached
colloquial lexeme with three closely related functions. First, it may serve as
a placeholder noun meaning roughly \trsl{thingy}, \trsl{that thing}, or
\trsl{what-do-you-call-it}. Secondly, it may function as a dismissive or
evaluative noun meaning \trsl{nonsense}, \trsl{excuses}, or \trsl{all that
fuss}. Thirdly, it may appear at the edge of an utterance as a playful hedge,
roughly \trsl{not really}, \trsl{just kidding}, or \trsl{I am only saying it
lightly}.

The reduced variant \ird{žma} shares these meanings, while its distribution
is associated with rapid, playful, or dismissive casual speech. This fixed
colloquial reduction is lexically specified; it does not extend the productive
deletion of an unstressed stem-final-syllable \ird{e} described in §\,\ref{sec:morphophon}.

\subsubsection{Category and core values}

The basic grammatical point is that \ird{žema} behaves as a noun-like base in
ordinary Iridian syntax. In its placeholder use it fills a nominal slot without
contributing a stable lexical class of its own:

\pex
\a
\begingl
\gla Tak žema Janku koupnik.//
\glb this thingy Janek.\Trs{} buy-\Trs{}-\Pf{}//
\glft \trsl{Janek bought that thing / that whatchamacallit.}//
\endgl
\a
\begingl
\gla Tak žema.//
\glb this nonsense//
\glft \trsl{That is nonsense / that is just an excuse.}//
\endgl
\xe

In the first example, \ird{žema} is simply a placeholder noun occupying the
patient slot of an ordinary transitive clause. In the second, the same form is
interpreted dismissively: the speaker rejects what has just been said or done as
empty, unserious, or merely performative. The lexical content is therefore weak
in both uses, but the syntactic behaviour remains nominal.

\subsubsection{Plurality}

The nominal behaviour of the item becomes still clearer when it appears in
plural form. Like other Class III nouns, \ird{žema} takes the plural direct
\ird{žemy} (syncretic with the singular transitive), the plural transitive
\ird{žem} (zero-ending on the consonant stem), and the plural indirect
\ird{žemám}. In this plural configuration \ird{žema} denotes a loose aggregate:
\trsl{all those things}, \trsl{all that stuff}, or, in a more evaluative
context, \trsl{all that nonsense}.

\ex
\begingl
\gla Žemy Marek koupnik.//
\glb thingy.\Dir{}.\Pl{} Marek.\Trs{} buy-\Pf{}//
\glft \trsl{Marek bought all that stuff.}//
\endgl
\xe

\subsubsection{Clause-peripheral hedge use}

The same colloquial base may also appear clause-peripherally, especially in the
reduced form \ird{žma}. In this use it no longer behaves like an ordinary noun
inside the clause, but as a parenthetical hedge by which the speaker playfully
withdraws, softens, or reframes what has just been said. Even so, it should not
be treated as a member of the formal particle hierarchy. Its contribution is
interactional and parenthetical rather than structural in the narrow sense.

\pex
\begingl
\gla Janek v ščenek, žma.//
\glb Janek.\Dir{} already arrive-\Pf{} hedge//
\glft \trsl{Janek has already arrived, not really / just kidding.}//
\endgl
\xe

This utterance-edge use follows naturally from the same semantic bleaching seen
in the nominal uses. The speaker introduces a proposition and then immediately
marks it as playful, non-committal, or not to be taken at full face value.
Synchronically, then, \ird{žema} and \ird{žma} form a small colloquial system:
an underspecified noun in the syntax, an evaluative noun in stance-marking
predication, and a parenthetical hedge in clause-peripheral discourse use.

\subsubsection{Verbal use and productive
recategorisation}\label{sec:spoken-zema-verbal}

The noun-like base also enters verbal morphology as a colloquial stem,
whose vague lexical content permits a range of contextual interpretations.

The colloquial verbal stem is \ird{žemat-}, assigned to thematic class~I
(thematic consonant \ird{-t}), with antipassive stem \ird{žemač-}. Like
other class~I stems, \ird{žemat-} takes the valency and
aspect-alignment morphology described in Chapter~\ref{ch:verbs}, although its
lexical content remains as indeterminate as that of the nominal base.

In its actor-focused (antipassive) use, the stem describes the agent as engaging
in vague, performative, evasive, or unserious activity. The core meaning is most
readily glossed as \trsl{to mess around}, \trsl{to make excuses}, \trsl{to put
on a show}, or \trsl{to fuss about without achieving anything}, though the exact
nuance depends on context in the same way as the nominal base does. The
antipassive perfective \ird{žemaček} reports a completed episode of such
activity; the antipassive retrospective \ird{žemačac} presents it with
continuing present relevance and a characteristically evaluative charge---the
speaker frames the prior activity as having amounted to nothing or as having
been mere performance; the continuous \ird{žemata} describes it as an ongoing
state or habitual disposition:

\pex
\a
\begingl
\gla Anka žemaček.//
\glb Anka.\Dir{} žemat-\Antip{}-\Pf{}//
\glft \trsl{Anka was just messing around / making excuses.}//
\endgl
\a
\begingl
\gla Anka žemačac.//
\glb Anka.\Dir{} žemat-\Antip{}-\Ret{}//
\glft \trsl{Anka has been making nothing but excuses.}//
\endgl
\a
\begingl
\gla Anka žemata.//
\glb Anka.\Dir{} žemat-\Cont{}//
\glft \trsl{Anka is always fussing / is the type to make excuses.}//
\endgl
\xe

In its patient-focused (ergative) use, the stem \ird{žemat-} takes the
transitive marker \ird{-n-} and acquires a closely related but distinct reading:
\trsl{to handle carelessly}, \trsl{to give only a vague or makeshift treatment
to}, \trsl{to fudge}, or \trsl{to treat as a trivial matter}. In the ergative
aspects the agent appears in the transitive case and the patient in the
directional:

\pex
\a
\begingl
\gla Anka tak žematnik.//
\glb Anka.\Trs{} that.\Dir{} žemat-\Trs{}-\Pf{}//
\glft \trsl{Anka fudged that / handled that carelessly.}//
\endgl
\a
\begingl
\gla Anka tak žematnac.//
\glb Anka.\Trs{} that.\Dir{} žemat-\Trs{}-\Ret{}//
\glft \trsl{Anka has been treating that carelessly the whole time.}//
\endgl
\xe

Semantically, all verbal uses share a stable core even as contextual readings
vary widely. The predicate never names a concrete or specific action; it instead
attributes to the agent a \emph{stance}---vagueness, evasiveness, theatricality,
or dismissal---and that stance is precisely what the nominal base also
contributes. \ird{Žemaček} in one context means \trsl{she was just messing
around}; in another, \trsl{he was making excuses}; in a third, \trsl{they were
fussing over trivial extras}. What remains constant is the evaluative frame: the
speaker characterises the activity not by what was done but by how it was done,
and the \textit{how} is always some variant of unserious, unfocused, or merely
performative.

The thematic alternation \ird{žemat-}/\ird{žemač-} and the aspect forms
illustrated above follow the ordinary class~I pattern. The construction's
colloquial character resides in the stem's vague, evaluative meanings and
register, while its inflection determines alignment and viewpoint in the
usual way.

\section{The \ird{ty}-supplement in spoken interaction}
\label{sec:spoken-ty}

The appositional particle \ird{ty} and the supplement construction it heads
are introduced in the nominal chapter (\S\,\ref{sec:nm-appos}). The present
section treats the specifically interactional properties of the
right-peripheral variant, which is substantially more productive in spoken
than in written language and which serves a distinct set of discourse functions
in everyday speech.

\subsection{Referent-first ordering as a spoken strategy}
\label{sec:spoken-ty-strategy}

Spoken language imposes real-time production constraints that written language
does not. A speaker who launches a noun phrase cannot always plan its full
modifier structure before the head noun must be uttered; equally, elaborating
descriptions may become relevant only as the discourse develops, after the
referent has already been introduced. The right-peripheral \ird{ty}-phrase
directly serves both situations. The speaker commits to a head noun,
establishes the referent, and then appends the clarification or specification
via \ird{ty} as a subsequent informational unit:

\pex
\a \ird{pres ty že vednikový i}
   \trsl{the man---that is, the one I saw}
\a \ird{nuž ty mrocná i}
   \trsl{the city---the beautiful one}
\a \ird{Karel ty mlazka}
   \trsl{Karel---my brother, you know}
\xe

\noindent This ordering reflects the referential logic of spoken discourse: the
head noun anchors the referent in the common ground and the \ird{ty}-phrase
adds identifying or descriptive content that the hearer can use to resolve or
enrich the reference. The construction is particularly valuable when the
describing material is long or syntactically complex, since a heavy prenominal
modifier stack---a nominalised relative clause, a multi-word adjectival phrase,
a postpositional elaboration---requires planning the entire modifier before
uttering the head noun. Right-peripheral \ird{ty} makes it possible to deliver
the head noun immediately and defer the specifying material.

\subsection{Repair and afterthought uses}
\label{sec:spoken-ty-repair}

The most distinctively spoken use of the \ird{ty}-supplement is as a repair
and afterthought device. A speaker may introduce a referent, continue the
utterance, and then append a \ird{ty}-phrase retroactively to specify it more
precisely:

\pex
\ird{pres vednik, ty dumě na i}
   \trsl{I saw the man---the one at the house}
\xe

% Add how to use with Not for disambiguation

In such cases the \ird{ty}-phrase is structurally detached from the preceding
clause or \synhead{NP}. It is not embedded inside the \synhead{NP} it
elaborates; it is appended to the right edge of the utterance as a freestanding
supplement, anchored to the earlier referent anaphorically via \ird{i}. The
prosodic signature is a short break before \ird{ty}---a slight pause or the
falling-then-rising contour typical of repair and self-correction---followed by
the supplement at a somewhat lower pitch level.

This repair-like use extends to cases where the \ird{ty}-phrase follows a
complete finite clause and reaches back to a referent introduced earlier in the
utterance:

\pex
\a \ird{Vednik, ty dumě na i}
   \trsl{I saw [him]---the one at the house}
\xe

Here \ird{ty} is not repairing an error but elaborating a
discourse referent that was already recoverable from context (via zero anaphora
in the main clause), offering the hearer additional identifying information as
an aside.

As described in \S\,\ref{sec:nm-appos}, the anaphoric pronoun \ird{i}
appears inside the \ird{ty}-phrase whenever the supplement's descriptive content
is non-nominal and requires a pronominal head to attach to. In spoken use this
is the normal case: spoken \ird{ty}-supplements typically contain adjectival or
relative-clause material, and \ird{i} serves as the pronominal host that allows
this material to form a well-headed nominal phrase while observing
head-finality. The resulting pattern, \ird{ty} \textsc{[modifier i]}, is highly
productive and has a characteristic phrasal contour: \ird{ty} receives mild
prominence, the modifier retains ordinary lexical stress, and \ird{i} is
typically reduced or deaccented at the right edge of the supplement. The
relative prominence of these words does not alter the initial-stress rule
(§\,\ref{sec:stress}).

Right-peripheral \ird{ty}-phrases are unmarked in spoken interaction. Their
frequency, their use in repair contexts, and their role as a strategy for
deferring heavy modifier material all reflect production-driven motivations that
are native to spontaneous speech. In formal and written language the
construction remains available but is marked: a right-peripheral \ird{ty}-phrase
in writing carries a deliberate elaborative implication, foregrounding the
supplement as a rhetorically distinct informational act rather than presenting
it as an integrated modifier. Heavy adnominal material---nominalised relative
clauses that would be unwieldy prenominally---can motivate the construction even
in formal prose without special effect; a short adjective or a simple nominal
going postnominal in writing, by contrast, is interpretable as emphatic or
contrastive. The contrast with the unmarked prenominal strategy is always
available to the writer as a pragmatic resource.
